% Figure from MATH 493 notes, section 71. Compile with the notes preamble.
\begin{tikzpicture}[>=Stealth, font=\scriptsize]
\draw[->, black!55, thin] (-2.3,0) -- (2.5,0) node[right] {$x$};
\draw[->, black!55, thin] (0,-2.3) -- (0,2.5) node[above] {$y$};
\draw[black!40] (0,0) circle (1.6);
\draw[blue!70!black, very thick] (25:-2.2) -- (25:2.2) node[right] {$L = \mathbb{R}(a, b)$};
\draw[red!70!black, very thick] (115:-2.2) -- (115:2.2) node[above] {$JL = \mathbb{R}(-b, a)$};
\draw[->, thick, blue!70!black] (0,0) -- (25:1.6);
\draw[->, thick, red!70!black] (0,0) -- (115:1.6);
\draw[->, black!70] (25:0.7) arc (25:115:0.7) node[midway, above right=-2pt] {$90^\circ$};
\node[anchor=west, align=left] at (4.3,1.2) {$J(a, b) = (-b, a)$ is perpendicular to $(a, b)$};
\node[anchor=west, align=left] at (4.3,0.4) {so $J$ maps every line $L$ through $0$\\to the perpendicular line $JL \neq L$};
\node[anchor=west, align=left] at (4.3,-0.6) {no line is a subrepresentation:\\$\mathbb{R}^2$ is simple over $\mathbb{R}$};
\end{tikzpicture}
